\section{راهنمای تنظیم بهره‌ها}

در این بخش، تمام بهره‌های قابل تنظیم سامانه پاندول معکوس با چرخ واکنشی را از دیدگاه ریاضی و نظریه کنترل بررسی می‌کنیم. برای هر کنترلر، نقش هر بهره در دینامیک حلقه‌بسته، اثر افزایش یا کاهش آن و رویه عملی تنظیم را ارائه می‌دهیم.

\subsection{دینامیک سامانه}

سامانه مورد بررسی یک سامانه الکترومکانیکی کوپل با پنج حالت است:

\begin{equation}
\mathbf{x} = \begin{bmatrix} \theta \\ \dot{\theta} \\ \varphi \\ \dot{\varphi} \\ i_a \end{bmatrix}
\end{equation}

که در آن $\theta$ زاویه پاندول نسبت به وضعیت عمودی، $\varphi$ زاویه چرخ نسبت به بازوی پاندول و $i_a$ جریان آرمیچر موتور است.

\subsubsection{زیرسامانه مکانیکی}

دینامیک مکانیکی توسط ماتریس اینرسی کوپل $2\times2$ حاکم بر سامانه توصیف می‌شود:

\begin{equation}
\mathbf{M}\begin{bmatrix} \ddot{\theta} \\ \ddot{\varphi} \end{bmatrix} = \mathbf{f}
\end{equation}

\begin{equation}
\mathbf{M} = \begin{bmatrix} M_{11} & M_{12} \\ M_{12} & M_{22} \end{bmatrix}
\end{equation}

عناصر این ماتریس به صورت زیر تعریف می‌شوند:

\begin{align}
M_{11} &= I_p + m_w L^2 + I_{w,\text{eff}} \\
M_{12} &= I_{w,\text{eff}} \\
M_{22} &= I_{w,\text{eff}}
\end{align}

که در آن $I_{w,\text{eff}} = I_w + J_r N^2$ مجموع اینرسی چرخ و اینرسی بازتابیده روتور از طریق جعبه‌دنده است.

بردار نیروی محرکه شامل جمله‌های گرانشی، گشتاور موتور و اصطکاک است:

\begin{align}
f_1 &= (m_p \ell_c + m_w L)\,g\sin\theta - N K_t i_a - b\,\dot{\theta} + \tau_{\text{ext}} \\
f_2 &= N K_t i_a - b_{w,\text{eff}}\,\dot{\varphi}
\end{align}

که در آن $b_{w,\text{eff}} = b_w + N^2 b_{\text{visc}}$ ضریب میرایی مؤثر چرخ است. حل این دستگاه معادلات با قاعده کرامر شتاب‌های زاویه‌ای را به‌دست می‌دهد:

\begin{equation}
\ddot{\theta} = \frac{M_{22}\,f_1 - M_{12}\,f_2}{\Delta}, \qquad \ddot{\varphi} = \frac{M_{11}\,f_2 - M_{12}\,f_1}{\Delta}
\end{equation}

که $\Delta = M_{11}M_{22} - M_{12}^2$ دترمینان ماتریس اینرسی است.

\subsubsection{زیرسامانه الکتریکی}

مدار آرمیچر موتور با معادله دیفرانسیل مرتبه اول زیر مدل می‌شود:

\begin{equation}
L_a\,\frac{di_a}{dt} = V - R_a\,i_a - K_e N\,\dot{\varphi}
\end{equation}

جمله $K_e N \dot{\varphi}$ نیروی ضد محرکه الکتریکی (Back-EMF) است که کوپل بین دو حوزه الکتریکی و مکانیکی را ایجاد می‌کند. این جمله تضمین می‌کند که افزایش سرعت چرخ، جریان آرمیچر و در نتیجه گشتاور موتور را محدود می‌کند.

\subsubsection{اثر پارامترهای فیزیکی بر دینامیک}

جدول \ref{tab:physical-params} اثر هر پارامتر فیزیکی بر رفتار دینامیکی سامانه را خلاصه می‌کند.

\begin{table}[H]
\centering
\caption{اثر پارامترهای فیزیکی بر دینامیک سامانه}
\label{tab:physical-params}
\begin{tabular}{@{}lll@{}}
\toprule
پارامتر & نماد & اثر بر دینامیک \\
\midrule
جرم پاندول & $m_p$ & افزایش گشتاور گرانشی؛ دشوارتر شدن بالانس \\
طول پاندول & $L$ & افزایش گشتاور و اینرسی؛ کاهش بسامد طبیعی \\
اینرسی چرخ & $I_w$ & افزایش ظرفیت ذخیره تکانه زاویه‌ای؛ پاسخ نرم‌تر ولی کندتر \\
نسبت تبدیل (اتصال مستقیم) & $N=1$ & ضرب گشتاور در $N$؛ افزودن اینرسی روتور در $N^2$ \\
ثابت موتور & $K_t = K_e$ & گشتاور به ازای هر آمپر؛ اقتدار عملگر بالاتر \\
مقاومت آرمیچر & $R_a$ & محدودسازی جریان ماندگار: $i_{ss} = V/R_a$ \\
القای آرمیچر & $L_a$ & ثابت زمانی الکتریکی $\tau_e = L_a/R_a$؛ پاسخ کندتر جریان \\
میرایی مفصل & $b$ & اصطکاک محور پاندول؛ پایدارسازی حلقه‌باز \\
میرایی چرخ & $b_w$ & اصطکاک یاتاقان چرخ؛ اتلاف انرژی چرخ \\
ولتاژ بیشینه & $V_{\max}$ & حد اشباع عملگر؛ کران گشتاور قابل دستیابی \\
\bottomrule
\end{tabular}
\end{table}


\subsection{کنترلر PID بالانس}

\subsubsection{قانون کنترل}

کنترلر PID بالانس، ولتاژ فرمان را بر اساس زاویه پاندول، انتگرال خطا و مشتق زاویه محاسبه می‌کند:

\begin{equation}
V(t) = K_p\,\theta + K_i\int_0^t \theta\,d\tau + K_d\,\dot{\theta}
\label{eq:pid-law}
\end{equation}

خروجی کنترلر در بازه $[-V_{\max},\; V_{\max}]$ محدود می‌شود. این محدودسازی تضمین می‌کند که فرمان ولتاژ از محدوده مجاز عملگر فراتر نرود.

\subsubsection{نقش هر بهره}

جدول \ref{tab:pid-gains} نقش ریاضی و اثر افزایش هر بهره را نشان می‌دهد.

\begin{table}[H]
\centering
\caption{نقش بهره‌های کنترلر PID}
\label{tab:pid-gains}
\begin{tabular}{@{}llll@{}}
\toprule
بهره & نماد & نقش ریاضی & اثر افزایش \\
\midrule
تناسبی & $K_p$ & سختی: $K_p\theta$ & پاسخ سریع‌تر؛ نوسان در مقادیر بالا \\
انتگرالی & $K_i$ & حذف آفست ماندگار & حذف خطای DC؛ باد-up در مقادیر بالا \\
مشتقی & $K_d$ & میرایی: $K_d\dot{\theta}$ & کاهش اورشوت؛ حساسیت به نویز \\
\bottomrule
\end{tabular}
\end{table}

بهره $K_p$ نقش سختی معادل را ایفا می‌کند. افزایش آن، نیروی بازگرداننده را برای هر انحراف زاویه‌ای بیشتر می‌کند و در نتیجه سرعت پاسخ را افزایش می‌دهد. با این حال، اگر $K_p$ از حد بحرانی فراتر رود، سامانه وارد نوسان پایدار یا ناپایدار می‌شود.

بهره $K_d$ به عنوان میرایی مصنوعی عمل می‌کند. این بهره با ایجاد ولتاژ متناسب با سرعت زاویه‌ای، انرژی جنبشی پاندول را دفع می‌کند و از اورشوت جلوگیری می‌نماید. در مقابل، مقادیر بالای $K_d$ حساسیت به نویز اندازه‌گیری را افزایش می‌دهند.

بهره $K_i$ خطای حالت ماندگار را حذف می‌کند. این بهره به‌ویژه در حضور اغتشاش‌های ثابت یا عدم تقارن مکانیکی ضروری است. با این حال، اگر $K_i$ بیش از حد بزرگ باشد، پدیده باد-up انتگرالی رخ می‌دهد و نوسان‌های کندی ایجاد می‌شود.

\subsubsection{مکانیزم ضد باد-up}

برای جلوگیری از اشباع انتگرال، از روش توقف انتگرال‌گیری (Clamping) استفاده می‌شود. هنگامی که ولتاژ خروجی به حد اشباع برسد و انتگرال در جهت تشدید اشباع رشد کند، انتگرال‌گیری متوقف می‌شود:

\begin{equation}
\frac{d}{dt}\!\left(\int\theta\,d\tau\right) = \begin{cases} \theta & \text{اگر } |V| < V_{\max} \\ \theta & \text{اگر } |V| = V_{\max} \text{ و } \theta \cdot V < 0 \\ 0 & \text{در غیر این صورت} \end{cases}
\end{equation}

شرط $\theta \cdot V < 0$ تضمین می‌کند که انتگرال تنها زمانی متوقف می‌شود که رشد آن باعث تشدید اشباع شود. اگر خطا در جهت کاهش اشباع باشد، انتگرال‌گیری ادامه می‌یابد.

\subsubsection{رویه تنظیم}

\begin{enumerate}[label=\arabic*.]
    \item $K_i = 0$ و $K_d = 0$ قرار می‌دهیم. $K_p$ را به تدریج افزایش می‌دهیم تا پاندول اطراف وضعیت عمودی با بسامد قابل قبول نوسان کند.
    \item $K_d$ را اضافه می‌کنیم تا نوسان میرا شود. قاعده سرانگشتی:
    \begin{equation}
    K_d \approx \frac{2\sqrt{K_p \cdot I_{\text{eff}}}}{V_{\max}}
    \end{equation}
    که $I_{\text{eff}}$ اینرسی مؤثر دیده‌شده از محور پاندول است.
    \item $K_i$ کوچک (معمولاً $K_i < 0.1\,K_p$) برای حذف آفست باقیمانده اضافه می‌کنیم.
\end{enumerate}

نقطه شروع پیشنهادی برای پارامترهای فیزیکی پیش‌فرض: $K_p = 50$، $K_i = 0.1$، $K_d = 10$.


\subsection{کنترلر LQR بالانس}

\subsubsection{خطی‌سازی}

در اطراف نقطه تعادل عمودی ($\theta = 0$، $\dot\theta = 0$، $\dot\varphi = 0$، $i_a = 0$)، سامانه غیرخطی به فرم فضای حالت خطی زیر تبدیل می‌شود:

\begin{equation}
\dot{\mathbf{x}} = A\mathbf{x} + B\,V
\end{equation}

با بردار حالت $\mathbf{x} = [\theta,\;\dot\theta,\;\dot\varphi,\;i_a]^\top$. ماتریس‌های $A$ و $B$ از خطی‌سازی معادلات دینامیک به‌دست می‌آیند:

\begin{equation}
A = \begin{bmatrix}
0 & 1 & 0 & 0 \\
\frac{M_{22}\,G}{\Delta} & -\frac{M_{22}\,b}{\Delta} & \frac{M_{12}\,b_{w,\text{eff}}}{\Delta} & -\frac{(M_{22}+M_{12})\,N K_t}{\Delta} \\
-\frac{M_{12}\,G}{\Delta} & \frac{M_{12}\,b}{\Delta} & -\frac{M_{11}\,b_{w,\text{eff}}}{\Delta} & \frac{(M_{11}+M_{12})\,N K_t}{\Delta} \\
0 & 0 & -\frac{K_e N}{L_a} & -\frac{R_a}{L_a}
\end{bmatrix}
\end{equation}

\begin{equation}
B = \begin{bmatrix} 0 \\ 0 \\ 0 \\ 1/L_a \end{bmatrix}
\end{equation}

که در آن $G = (m_p \ell_c + m_w L)\,g$ ضریب گشتاور گرانشی و $\Delta = M_{11}M_{22} - M_{12}^2$ دترمینان ماتریس اینرسی است. سطر چهارم ماتریس $A$ دینامیک الکتریکی آرمیچر را نشان می‌دهد و ورودی کنترل $V$ مستقیماً بر مشتق جریان اثر می‌گذارد.

\subsubsection{قانون کنترل بهینه}

کنترلر LQR تابع هزینه درجه دوم زیر را کمینه می‌کند:

\begin{equation}
J = \int_0^\infty \left(\mathbf{x}^\top Q\,\mathbf{x} + V^\top R_{\text{lqr}}\,V\right) dt
\label{eq:lqr-cost}
\end{equation}

بازخورد بهینه به صورت زیر است:

\begin{equation}
V = -K\mathbf{x} = -R_{\text{lqr}}^{-1} B^\top P\,\mathbf{x}
\end{equation}

که $P$ جواب معادله جبری ریکاتی پیوسته است:

\begin{equation}
A^\top P + PA - PBR_{\text{lqr}}^{-1}B^\top P + Q = 0
\end{equation}

این معادله به صورت عددی حل می‌شود و ماتریس بهره $K$ را به‌دست می‌دهد. اگر حل ریکاتی همگرا نشود (مثلاً به دلیل پارامترهای فیزیکی نامعتبر)، کنترلر به PID بازمی‌گردد.

\subsubsection{تفسیر ماتریس وزن‌دهی}

ماتریس $Q$ به صورت قطری تعریف می‌شود:

\begin{equation}
Q = \text{diag}(q_\theta,\; q_{\dot\theta},\; q_{\dot\varphi},\; q_{i_a})
\end{equation}

جدول \ref{tab:lqr-weights} تفسیر هر وزن و اثر افزایش آن را نشان می‌دهد.

\begin{table}[H]
\centering
\caption{تفسیر وزن‌های ماتریس $Q$ در کنترلر LQR}
\label{tab:lqr-weights}
\begin{tabular}{@{}llll@{}}
\toprule
وزن & نماد & جریمه می‌کند & اثر افزایش \\
\midrule
$q_\theta$ & \texttt{lqr\_q\_theta} & انحراف زاویه از عمودی & تنظیم دقیق‌تر زاویه؛ اصلاح تهاجمی‌تر \\
$q_{\dot\theta}$ & \texttt{lqr\_q\_theta\_dot} & سرعت زاویه‌ای پاندول & میرایی بیشتر نوسانات \\
$q_{\dot\varphi}$ & \texttt{lqr\_q\_phi\_dot} & سرعت چرخ & محدودسازی چرخش چرخ \\
$q_{i_a}$ & \texttt{lqr\_q\_current} & جریان آرمیچر & کاهش تنش عملگر \\
$R_{\text{lqr}}$ & \texttt{lqr\_r} & تلاش کنترلی (ولتاژ) & کنترل محافظه‌کارانه‌تر \\
\bottomrule
\end{tabular}
\end{table}

نسبت $Q/R_{\text{lqr}}$ تعیین‌کننده میزان تهاجمی بودن کنترل است. نسبت بزرگ‌تر به معنای اولویت تنظیم حالت‌ها بر مصرف انرژی کنترلی و نسبت کوچک‌تر به معنای کنترل ملایم‌تر با مصرف کمتر است.

\subsubsection{رویه تنظیم}

\begin{enumerate}[label=\arabic*.]
    \item $R_{\text{lqr}} = 1$ و تمام وزن‌های $Q$ برابر ۱ قرار می‌دهیم.
    \item $q_\theta$ را برای تنظیم دقیق‌تر زاویه افزایش می‌دهیم (بازه پیشنهادی: ۱۰۰ تا ۵۰۰).
    \item در صورت تداوم نوسان، $q_{\dot\theta}$ را افزایش می‌دهیم (بازه پیشنهادی: ۱ تا ۱۰).
    \item در صورت چرخش بیش از حد چرخ، $q_{\dot\varphi}$ را افزایش می‌دهیم (بازه پیشنهادی: ۱۰ تا ۵۰).
    \item در صورت وجود پیک‌های جریان، $q_{i_a}$ را افزایش می‌دهیم (بازه پیشنهادی: ۰/۰۱ تا ۱).
    \item برای کاهش کلی تهاجم کنترل، $R_{\text{lqr}}$ را افزایش می‌دهیم.
\end{enumerate}

نقطه شروع پیشنهادی: $q_\theta = 100$، $q_{\dot\theta} = 1$، $q_{\dot\varphi} = 10$، $q_{i_a} = 0.01$، $R_{\text{lqr}} = 1$.


\subsection{کنترلر بالاآوردن مبتنی بر انرژی}

\subsubsection{انرژی پاندول}

انرژی مکانیکی پاندول (بدون احتساب چرخ) نسبت به حالت سکون عمودی تعریف می‌شود:

\begin{equation}
E_p = \frac{1}{2}I_p\,\dot\theta^2 + (m_p\ell_c + m_w L)\,g\,(\cos\theta - 1)
\label{eq:pendulum-energy}
\end{equation}

در وضعیت عمودی ساکن، $E_p = 0$ است. در هر وضعیت پایین‌تر، $E_p < 0$ خواهد بود. هدف کنترلر بالاآوردن، رساندن $E_p$ به صفر است.

\subsubsection{قانون پمپاژ انرژی}

کنترلر بالاآوردن از روش آستروم–فوروبو استفاده می‌کند:

\begin{equation}
V = -k_e\,(E_{\text{target}} - E_p)\,\dot\varphi
\label{eq:energy-pump}
\end{equation}

که $E_{\text{target}} = 0$ سطح انرژی متناظر با وضعیت عمودی و $k_e$ بهره بالاآوردن انرژی است. منطق علامت‌دهی به این صورت عمل می‌کند: هنگامی که $E_p < 0$ (زیر هدف)، عبارت $(E_{\text{target}} - E_p)$ مثبت است و ولتاژ متناسب با $-\dot\varphi$ تعیین می‌شود. این تضمین می‌کند که صرف‌نظر از جهت چرخش چرخ، انرژی به پاندول تزریق شود.

\subsubsection{اثر بهره‌ها}

\begin{table}[H]
\centering
\caption{پارامترهای کنترلر بالاآوردن انرژی}
\label{tab:energy-gains}
\begin{tabular}{@{}lll@{}}
\toprule
پارامتر & نماد & اثر \\
\midrule
$k_e$ & \texttt{energy\_swing\_up\_gain} & تهاجم پمپاژ؛ بیشتر $\leftarrow$ تزریق سریع‌تر انرژی، خطر عبور از هدف \\
$\dot\varphi_{\max}$ & \texttt{swing\_up\_max\_wheel\_speed} & محدودکننده سرعت چرخ؛ ولتاژ صفر در $|\dot\varphi| \geq \dot\varphi_{\max}$ \\
\bottomrule
\end{tabular}
\end{table}

\subsubsection{مکانیزم‌های ایمنی}

برای جلوگیری از چرخش بی‌رویه چرخ و ناپایداری در نزدیکی وضعیت عمودی، چندین لایه حفاظتی در نظر گرفته شده است:

\begin{itemize}
    \item \textbf{محدودکننده سرعت چرخ:} هنگامی که $|\dot\varphi| \geq \dot\varphi_{\max}$، ولتاژ صفر می‌شود.
    \item \textbf{کاهش خطی:} در بازه $0.6\,\dot\varphi_{\max} < |\dot\varphi| < \dot\varphi_{\max}$:
    \begin{equation}
    V \leftarrow V \cdot \frac{\dot\varphi_{\max} - |\dot\varphi|}{0.4\,\dot\varphi_{\max}}
    \end{equation}
    \item \textbf{قطع سخت:} اگر $|\dot\varphi| > 1.5\,\dot\varphi_{\max}$، ولتاژ به صفر تنظیم می‌شود.
    \item \textbf{میرایی عبور انرژی:} هنگامی که $E_p > 0$ (انرژی بیش از هدف)، $V = -0.5\,k_e\,\dot\varphi$ برای دفع انرژی اضافی چرخ اعمال می‌شود.
    \item \textbf{تحریک سرعت پایین:} اگر $|\dot\varphi| < 0.5$ rad/s، تحریک فاز-مبنا با دامنه $0.3\,V_{\max}$ در جهت $\text{sign}(\sin\theta)$ اضافه می‌شود تا از توقف سامانه جلوگیری شود.
\end{itemize}

\subsubsection{رویه تنظیم}

\begin{enumerate}[label=\arabic*.]
    \item با $k_e = 1$ شروع می‌کنیم و مسیر انرژی را مشاهده می‌کنیم.
    \item اگر بالاآوردن کند است، $k_e$ را افزایش می‌دهیم (تا حدود ۵).
    \item اگر پاندول از وضعیت عمودی عبور می‌کند یا چرخ مکرراً اشباع می‌شود، $k_e$ را کاهش می‌دهیم.
    \item $\dot\varphi_{\max}$ را بر اساس محدودیت‌های مکانیکی چرخ و موتور تنظیم می‌کنیم (پیش‌فرض: ۵۰ rad/s).
\end{enumerate}


\subsection{بالاآوردن با خطی‌سازی بازخورد جزئی}

\subsubsection{شتاب مطلوب}

در این روش، دینامیک پاندول به صورت یک نوسانگر غیرخطی میرا شکل‌دهی می‌شود:

\begin{equation}
\ddot\theta_{\text{des}} = -k_{\text{pfl},p}\,\sin\theta - k_{\text{pfl},d}\,\dot\theta
\label{eq:pfl-desired}
\end{equation}

این رابطه، دینامیک مطلوب پاندول را به صورت یک سامانه جرم-فنر-دمپر غیرخطی تعریف می‌کند که $\theta$ را به سمت صفر هدایت می‌نماید.

\subsubsection{خطی‌سازی بازخورد}

از معادلات کوپل دینامیک، شتاب مورد نیاز چرخ محاسبه می‌شود:

\begin{equation}
\ddot\varphi_{\text{req}} = \frac{\Gamma\sin\theta - M_{11}\,\ddot\theta_{\text{des}}}{M_{12}}
\end{equation}

سپس گشتاور موتور مورد نیاز از معادله چرخ استخراج می‌شود:

\begin{equation}
\tau_m = \frac{M_{12}\,\ddot\theta_{\text{des}} + M_{22}\,\ddot\varphi_{\text{req}}}{N}
\end{equation}

\subsubsection{نگاشت ولتاژ شبه‌ایستا}

ولتاژ فرمان از ترکیب دو جمله به‌دست می‌آید:

\begin{equation}
V = R_a\,\frac{\tau_m}{K_t} + K_e N\,\dot\varphi
\label{eq:pfl-voltage}
\end{equation}

جمله اول، جریان لازم برای ایجاد گشتاور مطلوب را تأمین می‌کند و جمله دوم، نیروی ضد محرکه الکتریکی را جبران می‌نماید.

\subsubsection{اثر بهره‌ها}

\begin{table}[H]
\centering
\caption{بهره‌های کنترلر PFL}
\label{tab:pfl-gains}
\begin{tabular}{@{}llll@{}}
\toprule
بهره & نماد & نقش & اثر افزایش \\
\midrule
$k_{\text{pfl},p}$ & \texttt{pfl\_kp} & شکل‌دهی تناسبی: $-k_p\sin\theta$ & همگرایی سریع‌تر؛ تقاضای بیشتر از چرخ \\
$k_{\text{pfl},d}$ & \texttt{pfl\_kd} & میرایی: $-k_d\dot\theta$ & کاهش اورشوت نزدیک عمودی؛ کندی بالاآوردن \\
\bottomrule
\end{tabular}
\end{table}

\subsubsection{اشباع انرژی-آگاه}

هنگامی که $E_p > 0$ (انرژی پاندول از هدف عمودی فراتر رفته)، کنترلر به میرایی خالص چرخ تغییر وضعیت می‌دهد:

\begin{equation}
V = -K_e N\,\dot\varphi
\end{equation}

این تغییر از چرخش مداوم چرخ جلوگیری می‌کند و شرایط را برای تحویل به کنترلر بالانس فراهم می‌سازد.

\subsubsection{رویه تنظیم}

\begin{enumerate}[label=\arabic*.]
    \item با $k_{\text{pfl},p} = 5$ و $k_{\text{pfl},d} = 2$ شروع می‌کنیم.
    \item برای بالاآوردن سریع‌تر، $k_{\text{pfl},p}$ را افزایش می‌دهیم.
    \item در صورت عبور پاندول از وضعیت عمودی، $k_{\text{pfl},d}$ را افزایش می‌دهیم.
    \item اگر سرعت چرخ مکرراً به حد اشباع می‌رسد، $k_{\text{pfl},p}$ را کاهش می‌دهیم.
\end{enumerate}


\subsection{بالاآوردن با ضربه در سرعت صفر}

\subsubsection{مکانیزم عملکرد}

این روش، نقاط عطف سرعت زاویه‌ای پاندول ($\dot\theta = 0$) را شناسایی می‌کند و در آن لحظات، پالس ولتاژ کامل به چرخ اعمال می‌نماید:

\begin{equation}
V = \text{sign}(\theta) \cdot V_{\max} \quad \text{به مدت } T_{\text{impulse}}
\end{equation}

جهت پالس به گونه‌ای انتخاب می‌شود که گشتاور واکنشی، پاندول را در نیم‌دوره برگشت به سمت وضعیت عمودی هل دهد.

\subsubsection{پارامترها}

\begin{table}[H]
\centering
\caption{پارامترهای روش ضربه در سرعت صفر}
\label{tab:zv-params}
\begin{tabular}{@{}lll@{}}
\toprule
پارامتر & نماد & اثر \\
\midrule
$T_{\text{impulse}}$ & \texttt{zero\_velocity\_impulse\_duration} & پهنای پالس؛ بیشتر $\leftarrow$ انرژی بیشتر، خطر چرخش بیش از حد \\
$k_{zv}$ & \texttt{zero\_velocity\_swing\_gain} & (رزرو برای مقیاس‌دهی تناسبی) \\
\bottomrule
\end{tabular}
\end{table}

\subsubsection{مکانیزم‌های ایمنی}

\begin{itemize}
    \item اگر $|\dot\varphi| > 0.9\,\dot\varphi_{\max}$، ضربه لغو می‌شود.
    \item اگر $|\theta| < 0.05$ rad (نزدیک عمودی)، ضربه اعمال نمی‌شود و کنترل بالانس مسئول است.
    \item ضربه تنها زمانی فعال می‌شود که $E_p < 0$ باشد (کمبود انرژی وجود داشته باشد).
\end{itemize}

\subsubsection{تنظیم}

$T_{\text{impulse}}$ را برای تجمع سریع‌تر انرژی افزایش می‌دهیم (پیش‌فرض: ۰/۰۵ ثانیه). این روش برای سامانه‌هایی با نسبت دنده بالا مناسب‌تر است، زیرا یک پالس کوتاه ولتاژ، تکانه زاویه‌ای قابل توجهی در چرخ ایجاد می‌کند.


\subsection{کنترلر مد لغزشی}

\subsubsection{سطح لغزش}

سطح لغزش به صورت ترکیب خطی از حالات سامانه تعریف می‌شود:

\begin{equation}
s = c_1\,\theta + c_2\,\dot\theta + c_3\,\dot\varphi
\label{eq:sliding-surface}
\end{equation}

سطح $s = 0$ رابطه دینامیکی مطلوب بین زاویه پاندول، سرعت زاویه‌ای و سرعت چرخ را تعیین می‌کند. روی این سطح، سامانه به یک منیفولد مرتبه اول محدود می‌شود.

\subsubsection{قانون کنترل با لایه مرزی}

برای جلوگیری از پدیده چترینگ، از تابع اشباع به جای تابع علامت استفاده می‌شود:

\begin{equation}
V = -K\,\text{sat}\!\left(\frac{s}{\Phi}\right) - \eta\,s
\label{eq:smc-law}
\end{equation}

تابع اشباع به صورت زیر تعریف می‌شود:

\begin{equation}
\text{sat}(z) = \begin{cases} z & |z| \leq 1 \\ \text{sign}(z) & |z| > 1 \end{cases}
\end{equation}

در داخل لایه مرزی ($|s| \leq \Phi$)، کنترل رفتار خطی دارد و در خارج آن، به کنترل سوییچینگ کامل تبدیل می‌شود.

\subsubsection{نقش بهره‌ها}

\begin{table}[H]
\centering
\caption{بهره‌های کنترلر مد لغزشی}
\label{tab:smc-gains}
\begin{tabular}{@{}llll@{}}
\toprule
بهره & نماد & نقش ریاضی & اثر افزایش \\
\midrule
$c_1$ & \texttt{smc\_c1} & وزن زاویه در سطح لغزش & اصلاح قوی‌تر خطای زاویه \\
$c_2$ & \texttt{smc\_c2} & وزن سرعت در سطح لغزش & میرایی بیشتر در دینامیک سطح \\
$c_3$ & \texttt{smc\_c3} & وزن سرعت چرخ در سطح & کوپل دینامیک چرخ به شرط لغزش \\
$K$ & \texttt{smc\_k} & بهره قانون رسیدن & همگرایی سریع‌تر به $s=0$؛ چترینگ در مقادیر بالا \\
$\eta$ & \texttt{smc\_eta} & میرایی تناسبی بر $s$ & هموارسازی نزدیک‌شدن به سطح \\
$\Phi$ & \texttt{smc\_boundary} & ضخامت لایه مرزی & بیشتر $\leftarrow$ نرم‌تر ولی کمتر دقیق \\
\bottomrule
\end{tabular}
\end{table}

\subsubsection{دینامیک روی سطح لغزش}

روی سطح ($s = 0$):

\begin{equation}
c_1\,\theta + c_2\,\dot\theta + c_3\,\dot\varphi = 0
\end{equation}

این رابطه سامانه را به یک منیفولد مرتبه اول مقید می‌کند. مقدار ویژه دینامیک تقلیل‌یافته برای زیرسامانه پاندول تقریباً $-c_1/c_2$ است (هنگامی که $c_3$ کوچک باشد). بنابراین، نسبت $c_1/c_2$ سرعت همگرایی زاویه پاندول به صفر را تعیین می‌کند.

\subsubsection{رویه تنظیم}

\begin{enumerate}[label=\arabic*.]
    \item ابتدا $c_3 = 0$ قرار می‌دهیم. نسبت $c_1/c_2$ را برای تعیین ثابت زمانی حلقه‌بسته مطلوب انتخاب می‌کنیم: $\tau \approx c_2/c_1$.
    \item $c_3$ کوچک (۰/۵ تا ۲) برای ورود سرعت چرخ به بازخورد اضافه می‌کنیم.
    \item $\Phi$ را برای محدودسازی چترینگ تنظیم می‌کنیم (شروع با ۰/۰۵). در صورت مشاهده لرزش عملگر، افزایش می‌دهیم.
    \item $K$ را برای رسیدن سریع‌تر به سطح افزایش می‌دهیم (شروع با ۲).
    \item $\eta$ (شروع با ۰/۵) برای میرایی اضافی روی سطح اضافه می‌کنیم.
\end{enumerate}

نقطه شروع پیشنهادی: $c_1 = 10$، $c_2 = 5$، $c_3 = 1$، $K = 2$، $\eta = 0.5$، $\Phi = 0.05$.


\subsection{آستانه تحویل بالانس}

\subsubsection{شرط سوئیچ}

پس از آنکه کنترلر بالاآوردن پاندول را به نزدیکی وضعیت عمودی رساند، کنترل به کنترلر بالانس (LQR یا PID) تحویل داده می‌شود. شرط تحویل:

\begin{equation}
|\theta| < \theta_{\text{th}} \quad \text{و} \quad |\dot\theta| < \dot\theta_{\text{th}}
\label{eq:handoff}
\end{equation}

هر دو شرط باید همزمان برقرار باشند تا تحویل انجام شود.

\subsubsection{پارامترها و تنظیم}

\begin{table}[H]
\centering
\caption{پارامترهای آستانه تحویل}
\label{tab:handoff-params}
\begin{tabular}{@{}lll@{}}
\toprule
پارامتر & نماد & اثر \\
\midrule
$\theta_{\text{th}}$ & \texttt{upright\_angle\_threshold} & پنجره زاویه‌ای تحویل؛ بزرگ‌تر $\leftarrow$ سوئیچ زودتر \\
$\dot\theta_{\text{th}}$ & \texttt{upright\_velocity\_threshold} & پنجره سرعت؛ بزرگ‌تر $\leftarrow$ اجازه تحویل در سرعت بالاتر \\
\bottomrule
\end{tabular}
\end{table}

بازه پیشنهادی: $\theta_{\text{th}} \in [0.15,\; 0.5]$ rad (تقریباً ۹ تا ۲۹ درجه) و $\dot\theta_{\text{th}} \in [0.5,\; 2.0]$ rad/s.

نکته مهم این است که کنترلر بالانس باید برای شرط اولیه $(\theta_{\text{th}},\; \dot\theta_{\text{th}})$ پایدار باشد. اگر پس از تحویل، پاندول از وضعیت عمودی خارج می‌شود، باید حوزه جذب کنترلر بالانس را بزرگ‌تر کرد (افزایش $K_p$ یا $q_\theta$ در LQR) یا $\theta_{\text{th}}$ را کاهش داد.


\subsection{تعامل بین بهره‌ها}

\subsubsection{کوپل از طریق اشباع ولتاژ}

تمام کنترلرها ولتاژی در بازه $[-V_{\max},\; V_{\max}]$ تولید می‌کنند. کران گشتاور مؤثر عملگر:

\begin{equation}
\tau_{\max} = N \cdot K_t \cdot \frac{V_{\max}}{R_a}
\end{equation}

اگر بهره‌ها گشتاوری بیش از $\tau_{\max}$ تقاضا کنند، سامانه وارد اشباع می‌شود و عملکرد degrade می‌کند. برای کنترلر PID باید بررسی شود:

\begin{equation}
K_p \cdot \theta_{\max} + K_d \cdot \dot\theta_{\max} < V_{\max}
\end{equation}

\subsubsection{مقیاس‌دهی نسبت دنده}

نسبت تبدیل $N$ در سه نقطه از معادلات ظاهر می‌شود:
\begin{itemize}
    \item گشتاور چرخ: $\tau_w = N K_t i_a$ (ضرب در $N$)
    \item اینرسی روتور: $J_r N^2$ (ضرب در $N^2$)
    \item نیروی ضد محرکه: $K_e N \dot\varphi$ (ضرب در $N$)
\end{itemize}

در اتصال مستقیم ($N=1$)، گشتاور موتور مستقیماً به چرخ منتقل می‌شود و اینرسی روتور بدون ضریب به چرخ اضافه می‌گردد. بهره‌هایی که برای یک مقدار $N$ تنظیم شده‌اند، مستقیماً به مقدار دیگر قابل انتقال نیستند.

\subsubsection{پهنای باند الکتریکی}

پهنای باند مدار آرمیچر:

\begin{equation}
\omega_e = \frac{R_a}{L_a}
\end{equation}

اگر حلقه کنترل تغییرات جریانی سریع‌تر از $\omega_e$ تقاضا کند، دینامیک الکتریکی به گلوگاه سامانه تبدیل می‌شود. برای پارامترهای پیش‌فرض ($R_a = 1.2\;\Omega$، $L_a = 0.5\;\text{mH}$)، $\omega_e = 2400$ rad/s که بالاتر از پهنای باند مکانیکی سامانه است.


\subsection{جدول خلاصه بهره‌ها}

جدول \ref{tab:gains-summary} خلاصه‌ای از تمام بهره‌های کلیدی و هدف اصلی تنظیم هر کنترلر را ارائه می‌دهد.

\begin{table}[H]
\centering
\caption{خلاصه بهره‌ها و اهداف تنظیم}
\label{tab:gains-summary}
\begin{tabular}{@{}lll@{}}
\toprule
کنترلر & بهره‌های کلیدی & هدف اصلی تنظیم \\
\midrule
PID & $K_p$، $K_i$، $K_d$ & بالانس نزدیک عمودی \\
LQR & $Q$ قطری، $R_{\text{lqr}}$ & بهینه‌سازی تنظیم حالت در برابر تلاش کنترلی \\
بالاآوردن انرژی & $k_e$، $\dot\varphi_{\max}$ & پمپاژ انرژی بدون چرخش بیش از حد \\
PFL & $k_{\text{pfl},p}$، $k_{\text{pfl},d}$ & شکل‌دهی شتاب پاندول به سمت عمودی \\
ضربه سرعت صفر & $T_{\text{impulse}}$ & زمان‌بندی و بزرگی ضربه در نقاط عطف \\
SMC & $c_1,c_2,c_3,K,\eta,\Phi$ & پایدارسازی مقاوم با چترینگ محدود \\
تحویل & $\theta_{\text{th}}$، $\dot\theta_{\text{th}}$ & انتقال نرم از بالاآوردن به بالانس \\
\bottomrule
\end{tabular}
\end{table}


\subsection{گردش کار پیشنهادی تنظیم}

ترتیب پیشنهادی برای تنظیم کامل سامانه به صورت زیر است:

\begin{enumerate}[label=\arabic*.]
    \item \textbf{تنظیم پارامترهای فیزیکی:} جرم، طول، اینرسی و ثابت‌های موتور را مطابق سامانه واقعی تنظیم می‌کنیم.
    \item \textbf{بررسی پایداری حلقه‌باز:} شبیه‌سازی را بدون کنترلر اجرا و زمان افتادن طبیعی پاندول را مشاهده می‌کنیم.
    \item \textbf{تنظیم PID:} ابتدا با ساده‌ترین کنترلر شروع می‌کنیم. بالانس پایدار در محدوده $\pm 10°$ اطراف عمودی را تأمین می‌کنیم.
    \item \textbf{تنظیم LQR:} نسبت $Q/R_{\text{lqr}}$ را برای عملکرد مطلوب تنظیم و نتیجه را با PID مقایسه می‌کنیم.
    \item \textbf{تنظیم بالاآوردن:} با روش انرژی شروع می‌کنیم و $k_e$ را تنظیم می‌نماییم. سپس PFL را برای همگرایی سریع‌تر آزمایش می‌کنیم.
    \item \textbf{تنظیم آستانه تحویل:} انتقال نرم بدون اورشوت گذرا را تأمین می‌کنیم.
    \item \textbf{افزودن اغتشاش:} مقاوم بودن سامانه را با پالس‌های گشتاور و نویز آزمایش می‌کنیم.
    \item \textbf{آزمایش SMC:} در صورت نیاز به مقاوم بودن در برابر عدم قطعیت پارامتری، کنترلر مد لغزشی را تنظیم می‌کنیم.
\end{enumerate}

این ترتیب تضمین می‌کند که هر مرحله بر پایه مرحله قبلی بنا شود و عیب‌یابی در هر گام امکان‌پذیر باشد.